\documentclass[../../../include/open-logic-section]{subfiles}

\begin{document}

\olfileid{sth}{z}{nat}
\olsection{اپنے قدرتی عدد چننے}

%Let us take it for granted that the informal principle \stagesinf{}
%can be justified. It is also, though, worth commenting on the
%particular axiom we extracted from that informal principle. 

\olref[infinity-again]{defnomega} وچ اسیں ’’قدرتی عدد‘‘ دی عبارت دی صاف
تعریف دتی۔ تسی ایس مقرر کیتی گل نوں کیویں سمجھو؟ ایہ کوئی وجودیات بارے
دعویٰ نہیں، بس کجھ سیٹاں نوں قدرتی عدد \emph{منّن} دا فیصلہ اے۔ ایس خیال
دیاں وجوہاں نوں اسیں \olref[sfr][rel][ref]{sec}، \olref[sfr][arith][ref]{sec}
تے \olref[sfr][infinite][dedekindsproof]{sec} وچ چھوہیا سی۔ پر اسیں
ایہ وجوہاں ہور وی صاف کر سکدے آں۔

ساڈا لامتناہیت دا مسلّمہ \citet{VonNeumann1925} دے کم مطابق اے۔
پر ہِیٹھلا ہور مسلّمہ اے، جیہڑا اسیں اوہدی تھاں اپنا سکدے ساں:

\begin{defish}
\emph{زرمَیلو دا \citeyear{Zermelo1908Untersuchungen} لامتناہیت دا مسلّمہ۔}
اک سیٹ $A$ اے جیہدے لئی $\emptyset \in A$ تے
$(\forall x \in A)\{x\} \in A$۔
\end{defish}

جے اسیں اپنے (فون نوئمان توں لئے) لامتناہیت دے مسلّمے دی تھاں زرمَیلو دا
مسلّمہ ورتیا ہوندا، تاں سانوں اویویں اک ڈیڈیکنڈ لامتناہی سیٹ تے ایس توں
اک ڈیڈیکنڈ الجبرا مل جاندا۔ زرمَیلو دے طریقے وچ وی ساڈے الجبرے دا خاص
رکن $\emptyset$ ہوندا (ساڈے $0$ دی تھاں)۔ پر یک بہ یک نقشہ
$x \mapsto \{x\}$ ہوندا، $x \mapsto x \cup \{x\}$ نہیں۔
ایہدا سبھ توں سدھا نتیجہ ایہ اے کہ زرمَیلو $2$ نوں $\{\{\emptyset\}\}$
منّدا اے، جدکہ اسیں (فون نوئمان دے نال) $2$ نوں $\{\emptyset, \{\emptyset\}\}$
منّ دے آں۔

لامتناہیت دے اک مسلّمے نوں دوجے دی تھاں کیوں چنیے؟ وڈی عملی وجہ ایہ اے
کہ فون نوئمان دا طریقہ لامتناہی درجیاں دے عدداں نوں سنبھالن لئی سوکھ نال
ودھایا جا سکدا اے۔ اسیں \olref[ordinals][]{chap} توں اگے ایہ گل ویکھاں گے۔
پر صرف \emph{حساب کرن} دے لحاظ نال دونیں طریقے اکّو جِنّا کم کرن گے۔
سو جے کوئی تہانوں آکھے کہ قدرتی عدد سیٹ \emph{ای نیں}، تاں ظاہر سوال اے:
\emph{اوہ کیہڑے سیٹ نیں؟}

ایس عین سوال نوں \citet{Benacerraf1965} نے مشہور کیتا۔ پر ایہ گل زور دے
کے کہنی چاہیدی اے کہ ایہ اوس گل دی بس سبھ توں مشہور مثال اے جیہڑی اسیں
پہلاں کئی واری ویکھ چکے آں۔ بنیادی گل ایہ اے کہ سیٹ نظریہ سانوں کئی
’’سمجھ آؤن والیاں‘‘ قسماں دیاں شےواں دی \emph{نقل بناؤن} دا طریقہ دیندا اے۔
حقیقی، ناطق، صحیح تے قدرتی عدد تاں نیں ای، پر مرتب جوڑیاں، فنکشن تے تعلق وی۔
پر سیٹ نظریہ سانوں نقل دا \emph{اکّو اک} طریقہ کدی نہیں دیندا۔
\emph{ہمیشہ} ہور طریقے ہوندے نیں جیہڑے سیدھی طرح ساڈے اوہناں ای کم آ سکدے سن۔

\end{document}
